\documentclass[german]{cv}

\hypersetup{
    pdftitle={Lebenslauf von Jeongjun Ahn},
    pdfauthor={Jeongjun Ahn}
}
\cvname{Jeongjun Ahn}
\cvlastupdated{Zuletzt aktualisiert im Juli 2026}

\begin{document}

    \placelastupdatedtext
    \begin{header}
        \textbf{\fontsize{24 pt}{24 pt}\selectfont Jeongjun Ahn}

        \vspace{0.3 cm}

        \normalsize
        \mbox{{\color{black}\footnotesize\faMapMarker*}\hspace*{0.13cm}Leipzig, Deutschland}%
        \kern 0.25 cm%
        \AND%
        \kern 0.25 cm%
        \mbox{\hrefWithoutArrow{mailto:tyyyu1268@gmail.com}{\color{black}{\footnotesize\faEnvelope[regular]}\hspace*{0.13cm}tyyyu1268@gmail.com}}%
        \kern 0.25 cm%
        \AND%
        \kern 0.25 cm%
        \mbox{\hrefWithoutArrow{tel:+49 152 31627865}{\color{black}{\footnotesize\faPhone*}\hspace*{0.13cm}0152 31627865}}%
        \kern 0.25 cm%
        \AND%
        \kern 0.25 cm%
        \mbox{\hrefWithoutArrow{https://linkedin.com/in/jeong-jun-ahn}{\color{black}{\footnotesize\faLinkedinIn}\hspace*{0.13cm}jeong-jun-ahn}}%
        \kern 0.25 cm%
        \AND%
        \kern 0.25 cm%
        \mbox{\hrefWithoutArrow{https://github.com/prosk-sudo}{\color{black}{\footnotesize\faGithub}\hspace*{0.13cm}prosk-sudo}}%
    \end{header}

    \vspace{0.3 cm - 0.3 cm}


    \section{Ausbildung}

        \begin{twocolentry}{
        \textit{Okt. 2023 – Mai 2026}}
            \textbf{Lancaster University Leipzig}

            \textit{BSc in Computer Science}
        \end{twocolentry}

    \section{Sprachen}

        \begin{onecolentry}
            \textbf{Koreanisch} (Muttersprache) \textperiodcentered\ \textbf{Englisch} (fließend) \textperiodcentered\ \textbf{Deutsch} (Grundkenntnisse)
        \end{onecolentry}

    \section{Projekte}

        \begin{twocolentry}{
        \textit{März 2026}}
            \textbf{DOMtutor - DOMjudge v9 Migration}

            \textit{Bachelorarbeit \textperiodcentered\ \href{https://github.com/DOMtutor}{github.com/DOMtutor}}
        \end{twocolentry}

        \vspace{0.10 cm}
        \begin{onecolentry}
            \begin{highlights}
                \item DOMtutor, eine Sammlung von Automatisierungstools rund um das Open-Source-Judging-System DOMjudge für Programmierwettbewerbe, mit einer fork-basierten, inkrementellen Strategie von DOMjudge v7.3.5 auf v9.0.0 migriert - über 16 Releases mit Breaking Changes hinweg.
                \item Rohe SQL-Abfragen und Datenmodelle der pyjudge-Bibliothek an ein grundlegend reorganisiertes Datenbankschema angepasst - restrukturierte Ablage von Executables, umbenannte und entfernte Spalten, neue Pflichtfelder und externe Identifier sowie eine neue task-basierte Judging-Pipeline.
                \item Dockerisierte Umgebung mit DOMjudge v7 und v9 im Parallelbetrieb für direkten Vergleich und Debugging aufgebaut sowie die problemtools-Integration im Kattis-Format für Problemverifikation und PDF-Generierung aktualisiert.
                \item Alle Kern-Workflows - Upload von Einstellungen, Problemen und Contests sowie das Judging von Einreichungen - Ende-zu-Ende gegen eine laufende DOMjudge-v9-Instanz verifiziert.
            \end{highlights}
        \end{onecolentry}

        \vspace{0.13 cm}

        \begin{twocolentry}{
        \textit{2025}}
            \textbf{Lancaster Under Attack}

            \textit{\href{https://github.com/prosk-sudo/Lancaster-Under-Attack}{Lancaster-Under-Attack}}
        \end{twocolentry}

        \vspace{0.10 cm}
        \begin{onecolentry}
            \begin{highlights}
                \item Top-Down-Zombie-Survival-Shooter in Java mit dem libGDX-Framework als Gruppenprojekt im zweiten Studienjahr entwickelt, mit mehreren Leveln, mausgesteuertem Schießen sowie Missionen, Gegner auszuschalten und NPCs zu retten.
            \end{highlights}
        \end{onecolentry}

        \vspace{0.13 cm}

        \begin{twocolentry}{
        \textit{2026}}
            \textbf{codewars.nvim}

            \textit{\href{https://github.com/prosk-sudo/codewars.nvim}{codewars.nvim}}
        \end{twocolentry}

        \vspace{0.10 cm}
        \begin{onecolentry}
            \begin{highlights}
                \item Neovim-Plugin in Lua entwickelt, mit dem sich Codewars-Katas von codewars.com direkt im Editor durchsuchen und lösen lassen, inspiriert vom Plugin leetcode.nvim.
            \end{highlights}
        \end{onecolentry}

        \vspace{0.13 cm}

        \begin{twocolentry}{
        \textit{2026}}
            \textbf{linkedin-cli}

            \textit{\href{https://github.com/prosk-sudo/linkedin-cli}{linkedin-cli}}
        \end{twocolentry}

        \vspace{0.10 cm}
        \begin{onecolentry}
            \begin{highlights}
                \item Verwaistes Open-Source-Kommandozeilentool in Python für die LinkedIn API v2 (ursprünglich von tigillo) übernommen und dessen Authentifizierung auf das erforderliche OpenID Connect migriert.
            \end{highlights}
        \end{onecolentry}

        \vspace{0.13 cm}

        \begin{twocolentry}{
        \textit{2026}}
            \textbf{music-scores}

            \textit{\href{https://github.com/prosk-sudo/music-scores}{music-scores}}
        \end{twocolentry}

        \vspace{0.10 cm}
        \begin{onecolentry}
            \begin{highlights}
                \item Laufendes privates Projekt zum Transkribieren und Setzen von Noten mit LilyPond.
            \end{highlights}
        \end{onecolentry}

    \section{Engagement \& Aktivitäten}

        \begin{twocolentry}{
        \textit{Nov. 2023}}
            \textbf{IBM Qiskit Fall Fest}
        \end{twocolentry}

        \vspace{0.10 cm}
        \begin{onecolentry}
            \begin{highlights}
                \item An IBMs weltweiter Hackathon-Reihe zum Quantencomputing teilgenommen und mit Qiskit Konzepte des Quantencomputings praktisch erkundet.
            \end{highlights}
        \end{onecolentry}

        \vspace{0.13 cm}

        \begin{twocolentry}{
        \textit{Jan. 2024}}
            \textbf{Global Game Jam 2024}

            \textit{\href{https://github.com/LUComp/ggj2024}{LUComp/ggj2024}}
        \end{twocolentry}

        \vspace{0.10 cm}
        \begin{onecolentry}
            \begin{highlights}
                \item Gemeinsam mit einer Gruppe von Studierenden der Universität ein Spiel im 48-Stunden-Format des Events konzipiert, entwickelt und getestet, zum Thema \glqq Make Me Laugh\grqq.
            \end{highlights}
        \end{onecolentry}

        \vspace{0.13 cm}

        \begin{twocolentry}{
        \textit{Jan. 2025}}
            \textbf{Global Game Jam 2025}

            \textit{\href{https://github.com/LUComp/ggj2025}{LUComp/ggj2025}}
        \end{twocolentry}

        \vspace{0.10 cm}
        \begin{onecolentry}
            \begin{highlights}
                \item Gemeinsam mit einer Gruppe von Studierenden der Universität ein Spiel im 48-Stunden-Format des Events konzipiert, entwickelt und getestet, zum Thema \glqq Bubbles\grqq.
            \end{highlights}
        \end{onecolentry}

        \vspace{0.13 cm}

        \begin{twocolentry}{
        \textit{Okt. 2024 – Mai 2026}}
            \textbf{LUComp \& Robotics Club} - \textit{Mitglied}
        \end{twocolentry}

        \vspace{0.10 cm}
        \begin{onecolentry}
            \begin{highlights}
                \item Regelmäßig an Treffen und Veranstaltungen des Clubs teilgenommen und sich mit anderen Mitgliedern zu Computing- und Robotik-Projekten ausgetauscht.
            \end{highlights}
        \end{onecolentry}

        \vspace{0.13 cm}

        \begin{twocolentry}{
        \textit{Nov. 2024 – Mai 2026}}
            \textbf{Schachclub, Lancaster University Leipzig} - \textit{Vorsitz}
        \end{twocolentry}

        \vspace{0.10 cm}
        \begin{onecolentry}
            \begin{highlights}
                \item Den Schachclub der Universität organisiert und geleitet sowie regelmäßige Treffen und Partien für die Mitglieder koordiniert.
            \end{highlights}
        \end{onecolentry}

        \vspace{0.13 cm}

        \begin{twocolentry}{
        \textit{Okt. 2025 – Mai 2026}}
            \textbf{Stellvertretende Studierendenvertretung}
        \end{twocolentry}

        \vspace{0.10 cm}
        \begin{onecolentry}
            \begin{highlights}
                \item Als Bindeglied zwischen Studierenden und Universität vermittelt und dabei unterstützt, akademische und administrative Anliegen einzubringen und zu klären.
            \end{highlights}
        \end{onecolentry}

    \section{Kenntnisse}

        \begin{onecolentry}
            \textbf{Programmiersprachen:} Python, C, Java, JavaScript, Lua, SQL, HTML/CSS, Lilypond
        \end{onecolentry}

        \vspace{0.13 cm}

        \begin{onecolentry}
            \textbf{Frameworks \& Bibliotheken:} FastAPI, PyTorch, scikit-learn, Hugging Face
        \end{onecolentry}

        \vspace{0.13 cm}

        \begin{onecolentry}
            \textbf{Tools \& Plattformen:} Linux, Bash, Git, Docker
        \end{onecolentry}

\end{document}
